% !TeX root = ../动手学强化学习笔记.tex
% 由主文档通过 \input 载入；请编译主文档。
\setcounter{chapter}{10} % 保留原书章号。
\chapter{模型学习与规划}
\label{ch:rl-planning}

真实交互通常比计算昂贵。模型学习从真实经验估计动作后果，规划再用模型进行额外价值更新。其作用是重复利用已有信息，而不是凭空创造关于真实环境的新证据。

\section{经验模型与模型估计}

记 $N(s,a,s')$ 为观测转移次数，$N(s,a)=\sum_{s'}N(s,a,s')$。对已访问状态动作对，经验模型为
\[
\widehat P(s'\mid s,a)=\frac{N(s,a,s')}{N(s,a)},\qquad
\widehat r(s,a)=\frac{\text{该对全部观测奖励之和}}{N(s,a)}.
\]
若需保留奖励与后继的相关性，可存储联合样本并一同抽取，或者分别估计给定后继的奖励分布。未访问对不能直接除以零，应暂不用于规划，或明确引入先验模型。

分布模型给出各后继概率，样本生成模型只需生成可能结果。只记住最近一次后继适合确定性模型的简化实现；随机环境中这样做会忽略真实转移的不确定性。随着环境变化，旧计数也可能失效，需要遗忘机制与再探索。

\section{Dyna 框架}

Dyna-Q 把直接学习与规划放在同一循环。输入交互预算、规划次数 $m$、步长与探索规则；初始化 $Q$ 和空模型。每次先采取真实动作，收到转移后执行一次 Q-learning 更新并更新模型；再重复 $m$ 次，从已访问状态动作对中抽样，用模型生成奖励和后继，执行同形式的价值更新。输出为价值表和由它导出的策略。

真实样本修正模型和价值，模拟样本只修正价值。若模型精确，增加规划通常能更充分求解当前模型；若模型有偏，则会更快接近一个错误模型的解。因此规划次数与模型质量必须同时考察。

\begin{example}[把终点奖励向前传播]
确定性链 $s_0\to s_1\to s_2\to\text{终止}$ 的奖励为 $0,0,1$，取 $\gamma=1/2$、初值零、更新步长一。首次经历该链时，前两个即时更新仍可能为零，末步得到 $Q(s_2)=1$。之后模型按反向顺序规划两次，即得 $Q(s_1)=1/2$、$Q(s_0)=1/4$，无需再次真实走完整条链。
\end{example}

\section{优先扫掠与规划资源分配}

优先扫掠（prioritized sweeping）维护变化可能影响的前驱。某状态价值显著改变时，把其前驱状态动作对按预期备份变化大小加入优先队列，优先更新，再继续向前传播。确定性环境的前驱易于枚举；随机模型需根据转移概率考虑影响。优先级决定算力投向何处，不改变真实样本数。

全状态扫描覆盖全面但浪费在无关区域；轨迹采样集中于策略可能访问的区域；决策时规划则围绕当前状态临时搜索。选择哪种方式取决于模型访问成本、状态空间大小与决策时限。

\section{搜索与模型误差}

蒙特卡洛树搜索（Monte Carlo tree search，MCTS）围绕当前状态建立搜索树，反复执行选择、扩展、模拟和回传。选择规则兼顾已观察收益与访问不足的分支，模拟给出叶节点后的回报估计，回传更新路径统计。若节点访问次数为 $N(s)$、动作访问次数为 $N(s,a)$，一种树内选择得分是 $Q(s,a)+c\sqrt{\log N(s)/N(s,a)}$，未访问动作单独优先处理；$c>0$ 与回报尺度有关。

搜索结果依赖模拟器或学习模型的正确性。在迷宫中通道突然关闭时，旧模型仍可能预测可通行，再多搜索也无法发现未被真实数据揭示的变化。模型更新、重新探索和规划预算共同决定适应过程。模型误差与搜索不足是两类不同问题，增加计算只能直接缓解后者。

\section{习题与实验}
\begin{enumerate}
\item 某状态动作对先后观察后继 $u,u,v,u$，奖励为 $1,0,2,1$，求经验转移和期望奖励；说明单独采样后继与奖励会丢失什么信息。
\item 改变链式例子的规划顺序，比较传播到起点所需的备份次数。
\item 在同一迷宫比较不同规划步数，再在中途关闭通道。分别记录真实交互量、规划次数、运行时间和环境变化后的恢复过程。
\end{enumerate}




